\meetingheader
    {Mar 09, 2026}
    {}
    {\meetinglinks{
        \meetinglink{Smelly Tool Descriptions (Feb 2026)}{https://arxiv.org/abs/2602.14878}
        \meetinglink{CoALA: Cognitive Architectures for Language Agents}{https://arxiv.org/abs/2309.02427}
        \meetinglink{Self-Evolving Agents Survey (Fang+ 2025)}{https://arxiv.org/abs/2502.01248}
    }}
    {\begin{itemize}
        \item Pipeline fixes: dropped DeepSpeed for HF Trainer, fixed image padding, replaced GPT-4o judge with local Qwen3.5-27B-FP8, fixed batch size OOM, made training limits configurable
        \item Tool description design: restructured 20+ tools to Purpose/Guidelines/Limitations per smelly-descriptions ablation
        \item 5 agent episodes on FineVision5: Ep~1 matches baseline (MMStar 30.3, OCRBench 333); failure modes from aggressive filtering and budget reduction
    \end{itemize}}
    {}



%% ============================================================
%% SECTION: TOOL DESCRIPTION DESIGN
%% ============================================================

\section{Tool Description Design}

\subsection{How tool definitions reach the model}

% TODO: source citations for this subsection
In tool-calling LLM systems, tool definitions (name, description, parameter schema) are not passed through a side channel --- they are serialized into the model's context window as part of the system prompt. Each provider uses a model-specific format (Figure~\ref{fig:tool-formats}): the tool description is rendered by a chat template into literal prompt text that competes for attention with task instructions and conversation history. HuggingFace Transformers exposes this via the \texttt{tools} parameter of \texttt{apply\_chat\_template}, which passes tool JSON schemas through the model's Jinja2 template to produce the final prompt string.

This means tool description quality has a direct effect on agent behavior: a vague or noisy description wastes context tokens and increases the chance of incorrect tool selection or malformed arguments. The description is, in effect, the API contract between the tool developer and the LLM policy.

% Sources: HF chat_templating docs (https://huggingface.co/docs/transformers/en/chat_templating_writing),
%   HF tool use docs (https://huggingface.co/docs/transformers/main/en/chat_extras),
%   OpenAI Harmony format (https://developers.openai.com/cookbook/articles/openai-harmony),
%   Meta Llama 3.1 prompt format (https://www.llama.com/docs/model-cards-and-prompt-formats/llama3_1/)

\definecolor{anthropicbg}{HTML}{FFF0E6}
\definecolor{anthropicframe}{HTML}{D97706}
\definecolor{harmonybg}{HTML}{E8F5E9}
\definecolor{harmonyframe}{HTML}{388E3C}
\definecolor{qwenbg}{HTML}{E8F1FC}
\definecolor{qwenframe}{HTML}{2A74B9}
\definecolor{llamabg}{HTML}{F3E8FF}
\definecolor{llamaframe}{HTML}{7C3AED}

\begin{figure}[H]
\centering
\small
\caption{The same tool JSON schema rendered into prompt text by four providers. gpt-oss-120b, Qwen~3.5, and Llama~3.1 verified via \texttt{apply\_chat\_template}; Anthropic format from API documentation. The description field (highlighted) appears verbatim in each --- it is literal prompt text, not metadata.}
\label{fig:tool-formats}

\begin{tcolorbox}[enhanced, colback=anthropicbg, colframe=anthropicframe, boxrule=0.8pt, left=4pt, right=4pt, top=4pt, bottom=3pt, sharp corners,
  attach boxed title to top left={xshift=4pt, yshift=-2pt},
  boxed title style={empty, size=minimal},
  coltitle=anthropicframe, fonttitle=\sffamily\bfseries\large,
  title={Anthropic (Claude)}]
\ttfamily\scriptsize
<functions>\\
<function>\{"name": "vlm\_filter",\\
\quad "description": \colorbox{yellow!30}{"Remove rows whose images}\\
\quad \colorbox{yellow!30}{are corrupt, too small, or extreme aspect}\\
\quad \colorbox{yellow!30}{ratios."},\\
\quad "parameters": \{"dataset\_name":\\
\quad\quad \{"type": "string"\}\}\}</function>\\
</functions>
\end{tcolorbox}

\vspace{-4pt}

\begin{tcolorbox}[enhanced, colback=harmonybg, colframe=harmonyframe, boxrule=0.8pt, left=4pt, right=4pt, top=4pt, bottom=3pt, sharp corners,
  attach boxed title to top left={xshift=4pt, yshift=-2pt},
  boxed title style={empty, size=minimal},
  coltitle=harmonyframe, fonttitle=\sffamily\bfseries\large,
  title={gpt-oss-120b}]
\ttfamily\scriptsize
<|start|>developer<|message|>\\
\# Tools\\
\\
\#\# functions\\
\\
namespace functions \{\\
\\
// \colorbox{yellow!30}{Remove rows whose images are corrupt, too small,}\\
// \colorbox{yellow!30}{or extreme aspect ratios.}\\
type vlm\_filter = (\_: \{\\
dataset\_name: string,\\
\}) => any;\\
\\
\} // namespace functions<|end|>
\end{tcolorbox}

\vspace{-4pt}

\begin{tcolorbox}[enhanced, colback=qwenbg, colframe=qwenframe, boxrule=0.8pt, left=4pt, right=4pt, top=4pt, bottom=3pt, sharp corners,
  attach boxed title to top left={xshift=4pt, yshift=-2pt},
  boxed title style={empty, size=minimal},
  coltitle=qwenframe, fonttitle=\sffamily\bfseries\large,
  title={Qwen 3.5}]
\ttfamily\scriptsize
<|im\_start|>system\\
\# Tools\\
\\
You have access to the following functions:\\
\\
<tools>\\
\{"type": "function", "function": \{"name": "vlm\_filter",\\
"description": \colorbox{yellow!30}{"Remove rows whose images are corrupt,}\\
\colorbox{yellow!30}{too small, or extreme aspect ratios."},\\
"parameters": \{"dataset\_name": \{"type": "string"\}\}\}\}\\
</tools><|im\_end|>
\end{tcolorbox}

\vspace{-4pt}

\begin{tcolorbox}[enhanced, colback=llamabg, colframe=llamaframe, boxrule=0.8pt, left=4pt, right=4pt, top=4pt, bottom=3pt, sharp corners,
  attach boxed title to top left={xshift=4pt, yshift=-2pt},
  boxed title style={empty, size=minimal},
  coltitle=llamaframe, fonttitle=\sffamily\bfseries\large,
  title={Llama 3.1}]
\ttfamily\scriptsize
<|start\_header\_id|>user<|end\_header\_id|>\\
\\
Given the following functions, please respond with\\
a JSON for a function call ...\\
\\
\{"type": "function", "function": \{"name": "vlm\_filter",\\
"description": \colorbox{yellow!30}{"Remove rows whose images are corrupt,}\\
\colorbox{yellow!30}{too small, or extreme aspect ratios."},\\
"parameters": \{"dataset\_name": \{"type": "string"\}\}\}\}\\
\\
Filter the dataset<|eot\_id|>
\end{tcolorbox}

\end{figure}

\subsection{Smelly tool descriptions}

Hasan et al.~\cite{smellyToolDesc2026} study 856 MCP tool descriptions and find 97.1\% contain at least one ``smell'' --- a structural deficiency such as unclear purpose, missing constraints, or leaked implementation details. 56\% have an unclear purpose specifically.

They decompose tool descriptions into six components and ablate each. Augmenting all six improves task success by 5.85 pp (median), but ablations show that \textbf{Examples} can be removed with no statistically significant degradation, while \textbf{Purpose} is essential in every effective combination.

\subsection{Our format: Purpose / Guidelines / Limitations}

\definecolor{purposebg}{HTML}{FDECEA}
\definecolor{guidelinesbg}{HTML}{FFF9E6}
\definecolor{limitationsbg}{HTML}{E6F0FA}

We adopt the minimal effective set from~\cite{smellyToolDesc2026} and rewrote all 20+ tool descriptions: \colorbox{purposebg}{\textbf{Purpose}} (what the tool does + what it returns), \colorbox{guidelinesbg}{\textbf{Guidelines}} (when to use it + ordering constraints), \colorbox{limitationsbg}{\textbf{Limitations}} (failure modes + performance constraints). Below are three verbatim docstrings from the codebase showing this format applied to tools at different pipeline stages.

\medskip
\noindent\textbf{\texttt{vlm\_filter}} \textit{(preselection --- image structural quality)}
\begin{tcolorbox}[enhanced, colback=white, colframe=black!60, boxrule=0.6pt, left=6pt, right=6pt, top=0pt, bottom=0pt, sharp corners, boxsep=0pt]
\colorbox{purposebg}{\parbox{\dimexpr\linewidth-2\fboxsep}{%
\ttfamily\scriptsize\vspace{4pt}
\textbf{Purpose}: Remove rows whose images are corrupt, too
small (any side <200px), or have extreme aspect ratios
(>3:1). Returns kept/rejected counts and per-reason
breakdown.\vspace{4pt}
}}\\[-1pt]
\colorbox{guidelinesbg}{\parbox{\dimexpr\linewidth-2\fboxsep}{%
\ttfamily\scriptsize\vspace{4pt}
\textbf{Guidelines}: Run on each dataset after profile\_datasets
and before text-based filters. Only requires the
dataset\_id --- thresholds are fixed. Safe on all
dataset types.\vspace{4pt}
}}\\[-1pt]
\colorbox{limitationsbg}{\parbox{\dimexpr\linewidth-2\fboxsep}{%
\ttfamily\scriptsize\vspace{4pt}
\textbf{Limitations}: Requires loading full image data into
memory (calls ensure\_full internally). Slow on large
datasets with many shards. Does not assess image
content quality --- only structural properties.\vspace{4pt}
}}
\end{tcolorbox}

\medskip
\noindent\textbf{\texttt{submit\_finetune}} \textit{(postselection --- training)}
\begin{tcolorbox}[enhanced, colback=white, colframe=black!60, boxrule=0.6pt, left=6pt, right=6pt, top=0pt, bottom=0pt, sharp corners, boxsep=0pt]
\colorbox{purposebg}{\parbox{\dimexpr\linewidth-2\fboxsep}{%
\ttfamily\scriptsize\vspace{4pt}
\textbf{Purpose}: Finetune the target model on the curated dataset. Blocks until training completes. Returns a job\_id (e.g.\ ``ft-abc123'') that identifies the trained model for evaluation.\vspace{4pt}
}}\\[-1pt]
\colorbox{guidelinesbg}{\parbox{\dimexpr\linewidth-2\fboxsep}{%
\ttfamily\scriptsize\vspace{4pt}
\textbf{Guidelines}: Run after save\_to\_disk. Pass the dataset\_id of the mixed dataset (e.g.\ ``vqav2+captcha+cosyn\_chart''). All training hyperparameters (learning rate, epochs, batch size) are fixed by the task YAML --- you only control which data to train on. The returned job\_id is required by submit\_eval.\vspace{4pt}
}}\\[-1pt]
\colorbox{limitationsbg}{\parbox{\dimexpr\linewidth-2\fboxsep}{%
\ttfamily\scriptsize\vspace{4pt}
\textbf{Limitations}: Blocking call --- typically takes 10--30 minutes depending on dataset size. Training is capped at max\_steps from the task config. Cannot be cancelled once started. All hyperparameters are fixed --- the agent has no control over training configuration.\vspace{4pt}
}}
\end{tcolorbox}

\medskip
\noindent\textbf{\texttt{submit\_eval}} \textit{(postselection --- evaluation)}
\begin{tcolorbox}[enhanced, colback=white, colframe=black!60, boxrule=0.6pt, left=6pt, right=6pt, top=0pt, bottom=0pt, sharp corners, boxsep=0pt]
\colorbox{purposebg}{\parbox{\dimexpr\linewidth-2\fboxsep}{%
\ttfamily\scriptsize\vspace{4pt}
\textbf{Purpose}: Evaluate a finetuned model on VLM benchmarks configured in the task YAML (e.g.\ MMStar, OCRBench, MMMU, HallusionBench, MMVet, MathVista). Blocks until all benchmarks complete. Returns per-benchmark scores normalized to [0,1].\vspace{4pt}
}}\\[-1pt]
\colorbox{guidelinesbg}{\parbox{\dimexpr\linewidth-2\fboxsep}{%
\ttfamily\scriptsize\vspace{4pt}
\textbf{Guidelines}: Run after submit\_finetune. Pass model\_id set to the job\_id returned by submit\_finetune (e.g.\ ``ft-abc123'') --- do NOT pass a filesystem path. Benchmarks are configured automatically from the task YAML. Use the returned scores to inform the next episode's data curation strategy.\vspace{4pt}
}}\\[-1pt]
\colorbox{limitationsbg}{\parbox{\dimexpr\linewidth-2\fboxsep}{%
\ttfamily\scriptsize\vspace{4pt}
\textbf{Limitations}: Blocking call --- evaluation time scales with number of benchmarks (typically 15--45 minutes total). Benchmarks and evaluation protocol are fixed by the task config --- the agent cannot select which benchmarks to run. Requires a running vLLM inference server and judge model.\vspace{4pt}
}}
\end{tcolorbox}

\subsection{Parameter defaults as implicit guidance}

We also shifted \textbf{parameter defaults} to prevent silent data loss. For example, \texttt{quality\_filter} had \texttt{min\_answer\_words=2}, which silently filtered valid single-word VQA answers (``yes'', ``3'', ``red''). Defaults changed to 0 (disabled) with activation logic moved to Guidelines: ``Only raise these for caption or instruction datasets.'' Default values are part of the tool schema the model sees --- unsafe defaults are another form of smelly description.


%% ============================================================
%% SECTION 2: EXPERIMENTAL SETUP
%% ============================================================

\section{Experimental Setup}
\label{sec:setup}

\subsection{Data}

Five datasets from FineVision5, each containing 10k samples of image--text pairs (VQA, captioning, OCR, instruction-following). The agent's sample budget is 10k total --- it must select and mix from the available 50k. One dataset (captcha) is consistently auto-excluded by \texttt{vlm\_filter} due to structural image failures, leaving 4 effective source datasets.

\subsection{Fine-tuning}

LLaVA-1.5-7B full fine-tune on a single H200 (140GB). Vision tower frozen; only the language model and multimodal projector are trained.

Two fixes from the previous setup:
\begin{enumerate}
    \item \textbf{Dropped DeepSpeed for HF Trainer + Accelerate.} DeepSpeed ZeRO-3 added complexity without benefit on a single GPU. Switching to plain HF Trainer via \texttt{accelerate launch} simplified the stack and eliminated rank-cleanup exit code bugs. Liger kernel~\cite{hsu2024liger} provides fused RoPE/RMSNorm/SwiGLU/CrossEntropy on the Llama backbone for comparable memory savings.
    \item \textbf{Fixed image padding.} HF's \texttt{CLIPImageProcessor} does resize $\to$ center-crop, which loses edge content on non-square images. We replicate the original LLaVA \texttt{expand2square()} --- padding to square with CLIP image mean as background --- so that center-crop becomes a no-op on the spatial dimension. The data collator also now pads token sequences properly (pad token $=$ eos, labels $=$ $-100$) instead of relying on DeepSpeed's implicit handling.
\end{enumerate}

\begin{table}[H]
\centering
\small
\caption{Fine-tuning hyperparameters. All fixed by the task configuration --- the agent controls only which data is selected.}
\label{tab:finetune-config}
\begin{tabular}{ll}
\toprule
\textbf{Parameter} & \textbf{Value} \\
\midrule
Epochs & 1 \\
Batch size (per device) & 2 \\
Gradient accumulation & 2 (effective batch 8) \\
Learning rate & 2e-5 (cosine, warmup 3\%) \\
Precision & bf16 \\
Optimizer & AdamW (fused) \\
Max steps & 200 (cap) \\
Max sequence length & 2048 tokens \\
Step time & $\sim$2.8s/step \\
\bottomrule
\end{tabular}
\end{table}

\subsection{Evaluation}

Six benchmarks via VLMEvalKit, with Qwen3.5-27B-FP8 as the LLM judge (served locally via vLLM). Inference uses vLLM for the fine-tuned model.

\begin{table}[H]
\centering
\small
\caption{Eval sanity check: LLaVA-1.5-7B zero-shot scores. Baseline column is a mix of paper-reported and VLMEvalKit-reproduced values; see Note column. Local Qwen3.5 judge within $\sim$1 point of baselines.}
\label{tab:judge-sanity}
\begin{tabular}{lccll}
\toprule
\textbf{Benchmark} & \textbf{Baseline} & \textbf{Ours (Qwen3.5 judge)} & \textbf{Source} & \textbf{Note} \\
\midrule
MMStar & 30.3 & 30.4 & \cite{chen2024mmstar} & from paper \\
OCRBench & 317 & 317 & \cite{liu2023ocrbench} & VLMEvalKit run (paper: 297) \\
HallusionBench & 27.6 & 27.6 & \cite{guan2024hallusionbench} & VLMEvalKit composite \\
MMMU\_DEV\_VAL & 35.3 & 35.3 & \cite{yue2024mmmu} & lmms-eval (val split) \\
MMVet & 31.1 & 30.2 & \cite{yu2024mmvet} & from paper \\
MathVista\_MINI & 26.1 & 25.4 & \cite{lu2024mathvista} & VLMEvalKit run \\
\bottomrule
\end{tabular}
\end{table}

\noindent Judge-dependent benchmarks (MMVet, MathVista) differ by $<$1 point from published baselines --- acceptable for relative comparisons across curation runs. The reward signal used by the agent is a normalized composite of these benchmark scores.


%% ============================================================
%% SECTION 3: EPISODE RESULTS
%% ============================================================

\section{Episode Results}

With all fixes applied, the agent runs the full loop on the setup described in Section~\ref{sec:setup}.

\begin{table}[H]
\centering
\small
\caption{Benchmark scores across agent episodes. Baseline is LLaVA-1.5-7B zero-shot. Ep~1 matches/beats baseline on MMStar and OCRBench with only 10k curated samples. Ep~3--4 show failure modes from aggressive filtering and budget reduction. Ep~5 is the first run with restructured tool descriptions (Purpose/Guidelines/Limitations).}
\label{tab:episode-results}
\begin{tabular}{lcccccccc}
\toprule
& \textbf{Samples} & \textbf{MMStar} & \textbf{OCRBench} & \textbf{Hallusion} & \textbf{MMMU} & \textbf{MMVet} & \textbf{MathVista} & \textbf{Note} \\
\midrule
Baseline & --- & 30.3 & 317 & 27.6 & 35.3 & 31.1 & 26.1 & zero-shot \\
\midrule
Ep~1 & 9,998 & \textbf{30.3} & \textbf{333} & 46.9 & 3.6 & --- & --- & uniform mix, judge failed \\
Ep~2 & 9,997 & --- & --- & --- & --- & --- & --- & eval crashed \\
Ep~3 & 9,999 & 1.1 & 320 & 45.2 & 3.2 & --- & --- & quality\_filter aggressive \\
Ep~4 & 4,999 & 0.9 & 333 & 47.5 & --- & --- & --- & halved budget, eval crashed \\
Ep~5 & 10,000 & 1.2 & \textbf{328} & \textbf{46.2} & \textbf{25.8} & 0.5 & 0.2 & new docstrings, no quality\_filter \\
\bottomrule
\end{tabular}
\end{table}

\paragraph{Episode summary.}

\begin{table}[H]
\centering
\small
\caption{Episode behavior summary. Steps = total tool calls. Pipeline shows the tool sequence. Ep~5 is the first episode with restructured tool descriptions and seed memories from Ep~1--4.}
\label{tab:episode-summary}
\begin{tabular}{clcl}
\toprule
\textbf{Ep} & \textbf{Pipeline} & \textbf{Steps} & \textbf{Outcome} \\
\midrule
1 & profile $\to$ explore(4) $\to$ vlm\_filter(4) $\to$ quality\_filter(4) $\to$ mix $\to$ ft $\to$ eval & 22 & done (reward 0.293) \\
2 & profile $\to$ explore(5) $\to$ vlm\_filter(5) $\to$ quality\_filter(5) $\to$ mix $\to$ ft & 20 & eval crashed \\
3 & profile $\to$ vlm\_filter(5) $\to$ quality\_filter(4) $\to$ rollback(2) $\to$ mix $\to$ ft(2) $\to$ eval & 25 & done (reward 0.204) \\
4 & profile $\to$ vlm\_filter(5) $\to$ mix(10k) $\to$ ft(timeout) $\to$ mix(5k) $\to$ ft $\to$ eval & 18 & eval crashed \\
\midrule
5 & profile $\to$ vlm\_filter(5) $\to$ mix $\to$ save $\to$ ft $\to$ eval & \textbf{12} & eval OOM (scores recovered) \\
\bottomrule
\end{tabular}
\end{table}

\paragraph{Ep~1: baseline-competitive on first attempt.}
The agent's first curation run matched baseline MMStar (30.3 vs 30.3) and beat OCRBench (333 vs 317) using a uniform temperature-sampled mix across 4 datasets (captcha auto-excluded by \texttt{vlm\_filter}). This demonstrates the pipeline produces viable training data.

\paragraph{Ep~3: aggressive filtering destroys MMStar.}
The agent applied \texttt{quality\_filter} with \texttt{min\_answer\_words=2}, silently removing valid single-word VQA answers. After rollback and retry, MMStar collapsed to 1.07\% while OCRBench held (320). This failure motivated the default parameter changes described in Section~1.3.

\paragraph{Ep~4: budget sensitivity.}
After a finetune timeout (1808s on 10k samples), the agent halved to 5k. MMStar collapsed further to 0.93\% while HallusionBench was unaffected (47.5 aAcc). MMStar requires diverse general knowledge that needs the full 10k budget.

\paragraph{Ep~5: reasoning trace with restructured tool descriptions.}
Episode~5 is the first run after the tool description rewrite (Section~1.2) and with seed memories from Ep~1--4 injected into the system prompt. It completed the full pipeline in \textbf{12 steps with zero wasted actions}. Below is the verbatim reasoning trace from the episode event log (\texttt{events.jsonl}).

\subparagraph{Step 0 --- Strategic planning (line 3, T+0s).}
Before any tool call, the agent reviews all prior episodes from seed memory:

\begin{quote}\small\ttfamily
\textbf{Review of past episodes}\\
- Episode 1 (balanced 25\% each) achieved the highest MMStar score (\textasciitilde30\%) -- essentially the baseline. OCRBench was also decent (33.3). This suggests a uniform mix provides good general reasoning ability.\\
- Episodes 2--4 suffered from crashes or severe MMStar drops. Key issues observed:\\
\quad 1. \textbf{Over-aggressive quality\_filter} (especially in Episode 3) likely removed useful reasoning instructions, collapsing MMStar to \textasciitilde1\%.\\
\quad 2. \textbf{Dataset size reductions} (Episode 4 halved to 5k) dramatically harmed MMStar, confirming we need the full 10k budget.\\
\quad 3. \textbf{Zero usage of captcha} -- it was never sampled, yet OCRBench may benefit from OCR-focused data.
\end{ttfamily}\end{quote}

\noindent The agent then formulates a revised strategy with a key decision --- to \textbf{skip text-based filtering entirely}:

\begin{quote}\small\ttfamily
\textbf{Proposed revised strategy}\\
2. \textbf{Image filtering only}: Apply vlm\_filter to each dataset to drop corrupt/too-small images. \textbf{Do NOT run quality\_filter} (or keep all thresholds at 0) to avoid unintentionally discarding valuable instruction text.
\end{ttfamily}\end{quote}

\noindent This is a direct learned correction from Ep~3's failure: the agent correctly attributes MMStar collapse to text filtering removing useful reasoning instructions, not to image quality issues.

\subparagraph{Steps 1--6 --- Efficient filtering (lines 45--106, T+1s to T+528s).}
The agent executes \texttt{profile\_datasets} followed by \texttt{vlm\_filter} on all 5 datasets. At step~3, captcha returns 0/10{,}000 rows kept. The agent adapts immediately --- its original plan wanted 10--15\% captcha for OCRBench, but the data contradicts this:

\begin{quote}\small\ttfamily
Whoa! captcha all images filtered out (likely low-res). That means captcha may not be usable for visual instruction (maybe it's text captcha images small). We may still consider using it? But images are all too small; we cannot use them. So we should drop captcha from mix.
\end{quote}

\subparagraph{Step 7 --- Proportional reasoning (line 107, T+532s).}
The agent realizes it filtered 10k-row samples but needs full-dataset effective sizes for \texttt{compute\_mix\_ratio}. It derives proportional estimates:

\begin{quote}\small\ttfamily
We don't have full counts after filtering, but we know original dataset sizes from profile (num\_rows). We can estimate filtered proportion = kept/original\_rows * total rows.\\[4pt]
From profile: vqav2 total 82772 rows, kept proportion = 9970/10000=0.997, so filtered rows \ensuremath{\approx} 82772*0.997 \ensuremath{\approx} 82540.\\
captcha total 113062 rows, kept proportion =0, filtered rows =0.\\
cosyn\_chart total 116814 rows, kept proportion = 0.9942, filtered \ensuremath{\approx} 116134.\\
llava\_instruct total 157710 rows, kept proportion =0.998 \ensuremath{\Rightarrow} \ensuremath{\approx} 157395.\\
ocrvqa total 165746 rows, kept proportion =0.9829 \ensuremath{\Rightarrow} \ensuremath{\approx} 162804.
\end{ttfamily}\end{quote}

\noindent All five proportional estimates are arithmetically correct. The agent then reasons about temperature selection before settling on \texttt{temperature=5} (default).

\subparagraph{Step 8 --- Error recovery (line 111, T+533s).}
\texttt{compute\_mix\_ratio} fails because the agent passed 4 sizes for 5 datasets. It diagnoses and fixes in one step:

\begin{quote}\small\ttfamily
The tool expects effective\_sizes count to match number of active datasets (including those with 0 rows?). It sees 5 active datasets (including captcha). We gave 4 sizes. Need to provide size for captcha (0).
\end{quote}

\subparagraph{Steps 9--12 --- Mix, save, finetune, eval (lines 154--166, T+534s to T+1909s).}
The agent adjusts rounding (\emph{``add 2 to llava\_instruct to make 2615, total 10000''}), then executes the remaining pipeline with clean single-parameter calls:

\begin{itemize}
    \item \texttt{submit\_finetune(dataset\_id="vqav2+captcha+cosyn\_chart+llava\_instruct+ocrvqa")} --- no hyperparameters, correct dataset store name
    \item \texttt{submit\_eval(model\_id="ft-e95ca9d3")} --- correct job\_id format, no filesystem path
\end{itemize}

\noindent Finetune succeeded in 1226s (\textasciitilde20 min). Eval failed at runtime due to GPU OOM from the finetune process not fully releasing memory; scores were recovered manually from cached VLMEvalKit outputs.

\subparagraph{Analysis.}
Ep~5 demonstrates several capabilities enabled by seed memories and structured tool descriptions:
\begin{enumerate}
    \item \textbf{Cross-episode learning}: The agent correctly identifies that \texttt{quality\_filter} caused Ep~3's MMStar collapse and explicitly avoids it.
    \item \textbf{Data-driven adaptation}: When captcha is filtered to 0 rows, the agent abandons its plan to include it rather than forcing a workaround.
    \item \textbf{Quantitative reasoning}: Proportional effective-size estimation from sample-level filter results to full dataset scale.
    \item \textbf{Correct tool usage}: 12 steps, 0 retries, 0 hallucinated parameters, correct pipeline ordering (profile $\to$ vlm\_filter $\to$ compute\_mix\_ratio $\to$ vlm\_mix $\to$ save\_to\_disk $\to$ submit\_finetune $\to$ submit\_eval).
\end{enumerate}

\paragraph{Operational observations.}
\begin{itemize}
    \item Data ops (profile, filter, mix, save) complete in $<$10 min total. Training and eval dominate wall time ($\sim$20 min finetune + $\sim$15--45 min eval per episode).
    \item Captcha dataset always dropped: all images fail \texttt{vlm\_filter} structural checks across all 5 episodes.
    \item Zombie subprocess management is a critical operational concern: finetune and eval launch as decoupled subprocesses that survive agent crashes. Stale GPU processes caused 3 of 5 episode failures (Ep~2, Ep~4 eval crashes; Ep~5 eval OOM).
\end{itemize}


%% ============================================================
%% NEXT STEPS
%% ============================================================

\meetingtasks{
    \item \todo{Complete Ep~5 eval and update results table}
    \item \todo{Run 5+ episode stress test: does the agent improve curation strategy from eval feedback across episodes?}
    \item \todo{Implement episodic memory: store per-episode summaries (curation decisions $\to$ benchmark deltas) for cross-episode learning}
    \item \todo{Add CLIP score filter and SemDeDup to tool suite}
    \item \todo{Measure token cost per episode and identify context window pressure points}
    \item \todo{Read multi-agent literature (51 papers collected) and extract design decisions for sub-agent architecture}
}
